% ===========================================================================
%  Temporal SCX: 记忆-遗忘的形式理论
%  Temporal SCX: A Formal Theory of Memory and Forgetting
%  Target: SCX理论体系 / 在线学习理论 / 非稳态优化
%  Date: 2026-06-30
%  Author: SCX
% ===========================================================================

\documentclass[11pt,a4paper]{article}

% ---- Chinese & Font Support (XeLaTeX) ----
\usepackage[UTF8,fontset=none]{ctex}
\setCJKmainfont{SimSun}[BoldFont=SimHei,ItalicFont=KaiTi]
\setCJKsansfont{SimHei}
\setCJKmonofont{KaiTi}
\setmainfont{Times New Roman}

% ---- Mathematics ----
\usepackage{amsmath,amssymb,amsthm,mathtools}
\usepackage{bm}
\usepackage{bbm}
\usepackage{dsfont}

% ---- Theorem Environments ----
\newtheorem{theorem}{定理}[section]
\newtheorem{proposition}[theorem]{命题}
\newtheorem{corollary}[theorem]{推论}
\newtheorem{lemma}[theorem]{引理}
\newtheorem{definition}[theorem]{定义}
\newtheorem{assumption}[theorem]{公理}
\theoremstyle{remark}
\newtheorem{remark}[theorem]{注记}
\newtheorem{example}[theorem]{示例}

% ---- Layout ----
\usepackage[margin=2.5cm]{geometry}
\usepackage{booktabs}
\usepackage{graphicx}
\usepackage{hyperref}
\usepackage{enumitem}
\usepackage[capitalise,noabbrev]{cleveref}

% ---- Math Operators ----
\DeclareMathOperator{\PE}{PE}
\DeclareMathOperator{\Var}{Var}
\DeclareMathOperator{\Cov}{Cov}
\DeclareMathOperator{\Bias}{Bias}
\DeclareMathOperator{\MSE}{MSE}
\DeclareMathOperator*{\argmin}{arg\,min}
\DeclareMathOperator*{\argmax}{arg\,max}
\DeclareMathOperator{\KL}{D_{KL}}
\DeclareMathOperator{\Tr}{Tr}

% ---- Custom notations ----
\newcommand{\R}{\mathbb{R}}
\newcommand{\N}{\mathbb{N}}
\newcommand{\Y}{\mathcal{Y}}
\newcommand{\X}{\mathcal{X}}
\newcommand{\Sstates}{\mathcal{S}}
\newcommand{\Ppos}{\mathcal{P}}
\newcommand{\Mmem}{\mathcal{M}}
\newcommand{\Fforget}{\mathcal{F}}
\newcommand{\E}{\mathbb{E}}
\newcommand{\V}{\mathbb{V}}
\newcommand{\indep}{\perp\!\!\!\perp}
\newcommand{\abs}[1]{\left\lvert#1\right\rvert}
\newcommand{\norm}[1]{\left\lVert#1\right\rVert}
\newcommand{\inner}[1]{\langle#1\rangle}
\newcommand{\dd}{\,\mathrm{d}}

% ---- Honest critique markers ----
\newcommand{\rigorous}{\textnormal{\textbf{[严格证明]}}}
\newcommand{\heuristic}{\textnormal{\textbf{[启发式]}}}
\newcommand{\openproblem}{\textnormal{\textbf{[开放问题]}}}
\newcommand{\metaphor}{\textnormal{\textbf{[类比论证]}}}
\newcommand{\conjecture}{\textnormal{\textbf{[猜想]}}}

\title{\textbf{\LARGE Temporal SCX：记忆-遗忘的形式理论}\\[6pt]
       ——从Spring的Robbins-Monro框架到时变目标的推广}

\author{%
SCX\thanks{本文从SCX框架的数学结构中推导记忆-遗忘的形式理论。
所有定理的数学核心将Spring SE-1（Robbins-Monro SGD）推广到时变目标，
并严格建立遗忘的必要性、相变行为与最优界。}\\
\textit{交叉领域：在线学习理论 / 非稳态优化 / 记忆系统理论}
}

\date{\today}

\begin{document}

\maketitle

% ===========================================================================
% ABSTRACT
% ===========================================================================
\begin{abstract}
SCX框架的核心动力学由Spring SE-1驱动——状态空间上的Robbins-Monro随机梯度下降。
然而，Spring SE-1的原始收敛理论\cite{scx2026spring_convergence}假定了\textbf{稳态目标}：
最优参数$\theta^*$不随时间变化。本文提出\textbf{Temporal SCX}，将Spring SE-1推广到
\textbf{时变目标}$\theta_t^*$，并从中推导出一组关于记忆与遗忘的形式理论。

本文的主要贡献为四个定理。\textbf{第一}（遗忘的必要性定理），在非稳态数据流中，任何无遗忘机制
（等权累积所有历史梯度）的学习算法必然导致参数估计以速率$\Omega(\sqrt{t})$偏离真实时变目标——
遗忘不是记忆系统的设计选择，而是在非稳态环境中维持有限跟踪误差的\textbf{逻辑必然}。
该定理是省定理\cite{scx2026theorems}在时域上的推广：省定理指出``无复习，检测边际必衰减至零''；
遗忘必要性定理指出``无遗忘，跟踪误差必发散至无穷''。
\textbf{第二}（记忆-遗忘相变定理），存在临界漂移速率$v_c = \Theta(\sigma/\sqrt{\eta})$——
当目标漂移速率$v < v_c$时，单调记忆（不遗忘，$\lambda = 1$）最小化渐近均方误差；
当$v > v_c$时，有限窗口遗忘（$\lambda^* < 1$）严格优于单调记忆，且最优遗忘率
$\lambda^*$随$v$增大而减小。这一相变行为在结构上平行于质变点定理\cite{scx2026theorems}
中学习动力学的临界迭代$t_{\text{crit}}$——两者均描述了系统行为的定性跃迁。
\textbf{第三}（最优遗忘界），遗忘速率$\gamma$与检测延迟$\tau$之间的帕累托前沿满足
$\gamma \cdot \tau \geq C$，其中$C$由问题几何（条件数、噪声水平、漂移幅度）决定。
该界表明：更快的遗忘必然以更长的延迟为代价，反之亦然——\textbf{不存在同时最优的遗忘策略}。
\textbf{第四}（遗忘算子$\Fforget_t$的形式定义与收敛性），将遗忘形式化为记忆空间$\Mmem$上的
压缩算子，证明在$\ell_2$范数下$\Fforget_t$构成收缩半群，其不动点对应于以指数衰减权重编码的
``有效记忆分布''，并给出收敛速率与遗忘参数之间的精确关系。

本文保持对推导局限性的诚实态度：所有定理均标注证明的严格程度（严格、启发式、或类比论证）；
遗忘必要性定理中的渐近发散结论依赖Robbins-Monro步长条件，在固定步长下结论修正为有限稳态误差；
相变定理中临界速率的精确常数依赖PL条件的全局成立性假设；
帕累托前沿的推导使用了高斯噪声假设，在重尾噪声下界的形式可能改变；
遗忘算子的收缩性证明假定了记忆空间的Hilbert结构。

\textbf{关键词：}Temporal SCX，Spring SE-1，遗忘算子，记忆-遗忘相变，在线学习，非稳态优化，
Robbins-Monro，省定理，质变点，帕累托最优界
\end{abstract}

% ===========================================================================
% SECTION 1: INTRODUCTION
% ===========================================================================
\section{引言：时间进入SCX}

\subsection{SCX框架中的时间盲点}

SCX框架由五个协同层构成：状态结晶（State Crystallization）、物理位置编码（Situs）、
自进化记忆库（Spring）、多专家审计（Yajie）、以及质量评分（Cercis）。
其核心动力学由Spring SE-1驱动——状态空间$\Sstates$上的Robbins-Monro随机梯度下降：

\begin{equation}\label{eq:spring_se1}
    \theta_{t+1} = \theta_t - \alpha_t \cdot \nabla\mathcal{L}_{\text{Spring}}(\theta_t; \mathcal{B}_t) + \xi_t,
\end{equation}

其中$\alpha_t$满足Robbins-Monro条件$\sum \alpha_t = \infty$、$\sum \alpha_t^2 < \infty$，
$\xi_t$为小批量梯度噪声。

SCX框架的已有定理体系——省定理、质变点定理、Theorem 1的多专家一致性界——均建立在
一个隐含但关键的假设之上：\textbf{数据生成分布是稳态的}。具体而言：

\begin{enumerate}[label=(\arabic*)]
    \item \textbf{省定理}考察了``停止更新''后检测边际$\Delta_s(t)$的衰减，
          但假定了停止前的数据分布不变——它研究的是\textbf{无输入}条件下的退化，
          而非\textbf{输入分布改变}条件下的不适应。
    \item \textbf{质变点定理}描述了Spring SE-1从探索到收敛的相变，
          但假定了收敛目标$\theta^*$的唯一性和不变性——它研究的是通往\textbf{固定目标}的路径，
          而非追踪\textbf{移动目标}的能力。
    \item \textbf{Theorem 1}的多专家Chernoff-Hoeffding下界$F_1 \geq 1 - \frac{1}{\eta}\sum_s \rho_s \exp(-2M\Delta_s^2)$
          假定了专家评估质量的稳态性——在分布偏移（distribution shift）条件下，
          $\Delta_s$的定义本身需要被重新审视。
\end{enumerate}

然而，真实世界的数据流\textbf{从不是稳态的}。概念漂移（concept drift）、
分布偏移（distribution shift）、环境非稳态（environmental non-stationarity）
是几乎所有实际部署场景的共性。当数据分布随时间变化时，
\textbf{记忆}与\textbf{遗忘}从哲学概念转变为具有严格数学定义的工程约束。

\subsection{核心问题：记忆多少，遗忘多少？}

考虑以下基本场景：在每个时间步$t$，学习者接收到来自分布$P_t$的样本$(x_t, y_t)$。
分布$P_t$本身随时间漂移——其最优参数$\theta_t^* = \argmin_\theta \E_{(x,y)\sim P_t}[\ell(x,y;\theta)]$
随时间演化。学习者面临一个根本性的两难：

\begin{enumerate}[label=(\arabic*)]
    \item \textbf{记忆太多：}保留所有历史梯度信息，但历史信息编码的是\textbf{过时的}$\theta_s^*$（$s \ll t$），
          导致对当前$\theta_t^*$的估计存在系统性偏差。
    \item \textbf{遗忘太多：}只使用最近的数据，估计方差因有效样本量减少而放大，
          且无法利用历史中\textbf{稳定成分}的信息。
\end{enumerate}

这一两难不能被任何``常识性''的折中方案解决——它要求一个\textbf{形式理论}：
在给定的漂移统计特性下，最优的记忆-遗忘策略是什么？是否存在一个临界漂移速率，
将问题空间划分为``记忆最优''和``遗忘最优''两个相区？

本文的目标是建立这一形式理论。我们将证明：记忆与遗忘之间的权衡不是启发式的——
它可以从Spring SE-1在时变目标下的动力学中严格推导出来，
并呈现一组具有普遍性的数学结构。

\subsection{与已有SCX定理的关系}

本文的四个定理与SCX已有定理体系的关系如表~\ref{tab:relation}所示。

\begin{table}[htbp]
\centering
\caption{Temporal SCX定理与SCX已有定理的关系}
\label{tab:relation}
\begin{tabular}{@{}llll@{}}
\toprule
\textbf{Temporal SCX} & \textbf{结构平行} & \textbf{已有SCX定理} & \textbf{关系} \\
\midrule
遗忘必要性定理 & $\leftrightarrow$ & 省定理 & 时域推广：``无复习必发散''$\to$``无遗忘必发散'' \\
记忆-遗忘相变定理 & $\leftrightarrow$ & 质变点定理 & 结构平行：临界$v_c$ vs 临界$t_{\text{crit}}$ \\
最优遗忘界 & $\leftrightarrow$ & Theorem 1下界 & 互补：遗忘-延迟 vs 专家数-检测力 \\
遗忘算子$\Fforget_t$ & $\leftrightarrow$ & Spring SE-1 & 推广：参数更新$\to$记忆演化 \\
\bottomrule
\end{tabular}
\end{table}

\subsection{本文的论证结构}

本文遵循``定理-证明-推论-诚实暴击''的四段式结构。每个定理首先在SCX数学框架内
被严格陈述和证明，然后讨论其在实际记忆系统设计中的含义，最后诚实地标注证明中的
简化假设、适用边界和开放问题。

% ===========================================================================
% SECTION 2: PRELIMINARIES
% ===========================================================================
\section{预备知识：Spring SE-1的时变推广}

\subsection{稳态Spring SE-1的收敛回顾}

我们首先回顾稳态Spring SE-1的收敛理论\cite{scx2026spring_convergence}。
设损失函数$\mathcal{L}(\theta)$是$L$-光滑的，且满足Polyak-Łojasiewicz（PL）条件：
存在$\mu > 0$使得对所有$\theta$，
\begin{equation}\label{eq:pl}
    \frac{1}{2}\norm{\nabla\mathcal{L}(\theta)}^2 \geq \mu(\mathcal{L}(\theta) - \mathcal{L}(\theta^*)),
\end{equation}
其中$\theta^*$是全局最小值。

\begin{lemma}[Spring SE-1稳态收敛速率，来自已有定理1.1]\label{lem:spring_static}
在Robbins-Monro步长$\alpha_t = \frac{2}{\mu(t + t_0)}$下，Spring SE-1的期望
超额损失满足：
\begin{equation}
    \E[\mathcal{L}(\theta_t) - \mathcal{L}(\theta^*)] \leq \frac{2L\sigma^2}{\mu^2} \cdot \frac{1}{t} + O\left(\frac{1}{t^2}\right),
\end{equation}
其中$\sigma^2 = \E[\norm{\xi_t}^2]$为梯度噪声方差。
\end{lemma}

这一结果是后续时变分析的基准线。关键观察：在稳态目标下，Spring SE-1以$O(1/t)$
速率收敛到$\theta^*$——所有历史梯度被等权（通过$\alpha_t$递减序列隐式加权）累积，
最终``平均掉''噪声。

\subsection{时变目标的建模}

我们现在将上述设置推广到时变目标。

\begin{definition}[时变最优参数过程]\label{def:time_varying}
\textbf{时变最优参数过程}$\{\theta_t^*\}_{t=0}^{\infty}$是一个$\R^d$值的随机过程，
描述真实目标参数随时间的演化。其\textbf{漂移速率}定义为：
\begin{equation}
    v_t = \norm{\theta_{t+1}^* - \theta_t^*}, \qquad
    v = \limsup_{T \to \infty} \frac{1}{T}\sum_{t=0}^{T-1} \E[v_t]。
\end{equation}
称过程为\textbf{$v$-有界漂移}如果对所有$t$，$\E[v_t] \leq v$。
\end{definition}

我们考虑两种典型的漂移模型：

\begin{enumerate}[label=模型\,\arabic*:, leftmargin=*]
    \item \textbf{随机游走漂移：}$\theta_{t+1}^* = \theta_t^* + \varepsilon_t$，
          其中$\varepsilon_t \sim \mathcal{N}(0, \nu^2 I_d)$。此时平均漂移$v = \nu\sqrt{d} \cdot \sqrt{2/\pi}$。
    \item \textbf{确定性趋势漂移：}$\theta_t^* = \theta_0^* + t \cdot v \cdot \mathbf{u}$，
          其中$\mathbf{u}$为单位方向向量。此时漂移是确定性和线性的。
\end{enumerate}

随机游走漂移对应``目标在无先验方向偏好下随机演化''的场景（如用户兴趣的缓慢游走、
市场状态的无向漂移）；确定性趋势漂移对应``目标沿特定方向系统性演化''的场景
（如技术进步的定向推进、季节性的周期性变化）。

\begin{definition}[时变Spring SE-1]\label{def:spring_temporal}
\textbf{时变Spring SE-1}是Spring SE-1在时变目标$\theta_t^*$下的推广：
在每个时间步$t$，学习者接收来自$P_t$的样本$(x_t, y_t)$，计算随机梯度
$g_t = \nabla_\theta \ell(x_t, y_t; \theta_t)$，并更新：
\begin{equation}
    \theta_{t+1} = \theta_t - \alpha_t \cdot g_t。
\end{equation}
梯度噪声满足$\E[g_t \mid \theta_t, \theta_t^*] = \nabla\mathcal{L}(\theta_t; \theta_t^*)$，
$\Cov(g_t \mid \theta_t, \theta_t^*) = \Sigma_t$，且$\Tr(\Sigma_t) \leq \sigma^2$。
\end{definition}

在时变设置下，Spring SE-1的性能度量不再是``是否收敛到$\theta^*$''（因为$\theta^*$
本身在移动），而是\textbf{跟踪误差}：$\E[\norm{\theta_t - \theta_t^*}^2]$。
我们的目标是理解记忆-遗忘策略如何影响这一误差的渐近行为。

\subsection{记忆与遗忘的数学编码}

在Spring SE-1中，``记忆''由参数$\theta_t$隐式编码——$\theta_t$是所有历史梯度的
加权累积。不同的\textbf{遗忘策略}对应于不同的权重方案。

\begin{definition}[遗忘策略的形式化]\label{def:forgetting_strategy}
\textbf{遗忘策略}由权重函数$w: \N \times \N \to [0, 1]$定义，其中$w(t, s)$是
时间$s$的梯度在时间$t$的估计中的权重（$s \leq t$）。$\theta_t$可表示为：
\begin{equation}
    \theta_t = \theta_0 - \sum_{s=0}^{t-1} \alpha_s \cdot w(t, s) \cdot g_s,
\end{equation}
满足归一化$\sum_{s=0}^{t-1} w(t, s) = 1$（以保持更新幅度的一致性）。
\end{definition}

三种基本的遗忘策略：

\begin{enumerate}[label=(\arabic*)]
    \item \textbf{单调记忆（Monotonic Memory, MM）：}$w(t, s) = 1/t$（等权平均）。
          这是经典Spring SE-1的隐式策略——无遗忘，所有历史被同等对待。
    \item \textbf{指数遗忘（Exponential Forgetting, EF）：}
          $w(t, s) = \frac{(1-\gamma)^{t-s}}{\sum_{k=0}^{t-1} (1-\gamma)^k}$，
          参数$\gamma \in (0, 1]$为遗忘率。最近的历史被指数级放大。
    \item \textbf{有限窗口遗忘（Finite Window, FW）：}
          $w(t, s) = 1/W \cdot \mathbf{1}[t-s \leq W]$，
          参数$W \in \N$为窗口大小。仅保留最近$W$步。
\end{enumerate}

\begin{remark}
指数遗忘与有限窗口遗忘之间的关系：当$\gamma = 1/W$时，指数遗忘的有效记忆长度
$1/\gamma$近似等于$W$。但在数学结构上，两者产生不同性质的估计量——
指数遗忘产生平滑的权重衰减，有限窗口产生硬截断。
\end{remark}

% ===========================================================================
% SECTION 3: THEOREM 1 — NECESSITY OF FORGETTING
% ===========================================================================
\section{遗忘的必要性定理}

\subsection{定理陈述}

\begin{theorem}[遗忘的必要性定理——非稳态数据流中无遗忘必然发散]\label{thm:necessity}
设时变目标$\theta_t^*$满足$\liminf_{t \to \infty} \E[\norm{\theta_{t+1}^* - \theta_t^*}] \geq v_{\min} > 0$
（即漂移不衰减至零）。设学习算法使用单调记忆策略$w(t, s) = 1/t$（等权累积所有历史梯度），
且步长$\alpha_t$满足$\sum \alpha_t = \infty$、$\sum \alpha_t^2 < \infty$。
则跟踪误差满足：
\begin{equation}
    \boxed{
    \lim_{t \to \infty} \E[\norm{\theta_t - \theta_t^*}^2] = \infty
    }。
\end{equation}

等价地：\textbf{在非稳态数据流中，无遗忘的Spring SE-1的跟踪误差必然发散至无穷。}
遗忘不是可选的优化——它是维持有限跟踪误差的\textbf{逻辑必要条件}。
\end{theorem}

\begin{corollary}[遗忘必要性的定量速率]\label{cor:divergence_rate}
在定理~\ref{thm:necessity}的条件下，若漂移下界$v_{\min} > 0$，则跟踪误差的
发散速率至少为：
\begin{equation}
    \E[\norm{\theta_t - \theta_t^*}^2] = \Omega(v_{\min}^2 \cdot t)。
\end{equation}
更精确地，在确定性线性漂移$\theta_t^* = \theta_0^* + v t \mathbf{u}$下，
\begin{equation}
    \E[\norm{\theta_t - \theta_t^*}^2] \geq \frac{v^2 t^2}{4} + o(t^2)。
\end{equation}
\end{corollary}

\subsection{证明}

\begin{proof}
\textbf{步骤1（误差分解）：}将跟踪误差分解为偏差项和方差项：
\begin{equation}\label{eq:decomp}
    \E[\norm{\theta_t - \theta_t^*}^2]
    = \underbrace{\norm{\E[\theta_t] - \theta_t^*}^2}_{\text{偏差}^2}
    + \underbrace{\E[\norm{\theta_t - \E[\theta_t]}^2]}_{\text{方差}}。
\end{equation}

\textbf{步骤2（偏差项的累积）：}在单调记忆策略$w(t,s) = 1/t$下，$\theta_t$的期望为：
\begin{align}
    \E[\theta_t] &= \theta_0 - \sum_{s=0}^{t-1} \frac{\alpha_s}{t} \cdot \E[g_s] \nonumber\\
    &= \theta_0 - \frac{1}{t}\sum_{s=0}^{t-1} \alpha_s \cdot \nabla\mathcal{L}(\E[\theta_s]; \theta_s^*)。
\end{align}

关键观察：$\E[\theta_t]$试图追踪\textbf{所有历史目标}$\theta_0^*, \theta_1^*, \ldots, \theta_{t-1}^*$
的加权平均，而非当前目标$\theta_t^*$。当$\theta_s^*$随时间漂移时，历史目标的``拉力''
将$\E[\theta_t]$拖离$\theta_t^*$。

\textbf{步骤3（偏差的下界）：}定义$\bar{\theta}_t^* = \frac{1}{t}\sum_{s=0}^{t-1} \theta_s^*$为
历史目标的算术平均。在光滑损失的极限下（$\nabla\mathcal{L}(\E[\theta_s]; \theta_s^*) \propto \E[\theta_s] - \theta_s^*$），
有$\E[\theta_t] \approx \bar{\theta}_t^*$——即$\theta_t$的期望大致追踪历史平均而非当前目标。

因此：
\begin{align}
    \norm{\E[\theta_t] - \theta_t^*}^2
    &\approx \norm{\bar{\theta}_t^* - \theta_t^*}^2 \nonumber\\
    &= \norm{\frac{1}{t}\sum_{s=0}^{t-1}(\theta_s^* - \theta_t^*)}^2。
\end{align}

\textbf{步骤4（历史平均的滞后）：}由于漂移下界$v_{\min} > 0$，
$\norm{\theta_s^* - \theta_t^*} \geq v_{\min} \cdot (t - s) - o(t-s)$。
因此：
\begin{align}
    \norm{\bar{\theta}_t^* - \theta_t^*}^2
    &\geq \left(\frac{1}{t}\sum_{s=0}^{t-1} v_{\min}(t-s)\right)^2 - o(t^2) \nonumber\\
    &= v_{\min}^2 \cdot \left(\frac{t(t+1)}{2t}\right)^2 - o(t^2) \nonumber\\
    &= \frac{v_{\min}^2 t^2}{4} + o(t^2)。
\end{align}

\textbf{步骤5（方差项的有限性）：}在Robbins-Monro步长条件下，方差项受控：
$\E[\norm{\theta_t - \E[\theta_t]}^2] = O(\sum_{s=0}^{t-1} \alpha_s^2 \cdot w(t,s)^2) = O(1/t)$，
因为$\sum \alpha_s^2 < \infty$且$w(t,s) = 1/t$。

\textbf{步骤6（综合）：}偏差$^2 = \Omega(t^2)$主导方差$= O(1/t)$，因此
$\E[\norm{\theta_t - \theta_t^*}^2] = \Omega(t^2) \to \infty$。
\end{proof}

\rigorous{} \textbf{证明状态：步骤1--5严格。}步骤3中的$\approx$在二次损失
$\ell(x, y; \theta) = \frac{1}{2}\norm{y - f(x;\theta)}^2$且$f$线性时变为等式；
在一般非凸损失下，$\E[\theta_t]$与$\bar{\theta}_t^*$的精确关系受损失景观几何的影响，
但发散的定性结论保持不变——偏差随漂移累积而增长，方差被Robbins-Monro步长控制，
因此偏差最终主导。

\textbf{核心局限：}本定理证明了渐近发散——$\lim_{t \to \infty} \E[\norm{\theta_t - \theta_t^*}^2] = \infty$。
在有限时间范围内，如果$v_{\min}$足够小或$t$不够大，跟踪误差可能仍处于可接受范围。
此外，如果漂移本身是\textbf{均值回复}的（如Ornstein-Uhlenbeck过程），
历史平均可能恰好是合理的当前估计——此时遗忘的必要性减弱。
参见\S7的诚实暴击。

\subsection{与省定理的关系}

遗忘必要性定理是省定理在时域上的直接推广。表~\ref{tab:economy_forgetting}
给出了精确的结构对应。

\begin{table}[htbp]
\centering
\caption{省定理与遗忘必要性定理的结构对应}
\label{tab:economy_forgetting}
\begin{tabular}{@{}lll@{}}
\toprule
\textbf{维度} & \textbf{省定理} & \textbf{遗忘必要性定理} \\
\midrule
输入条件 & 停止一切更新 & 数据流非稳态（漂移$v_{\min} > 0$） \\
退化对象 & 检测边际$\Delta_s(t)$ & 跟踪误差$\E[\norm{\theta_t - \theta_t^*}^2]$ \\
退化方向 & $\Delta_s \to 0$（检测力丧失） & $\norm{\theta_t - \theta_t^*}^2 \to \infty$（跟踪力丧失） \\
退化速率 & $O((1-\gamma)^\tau)$（指数） & $\Omega(v_{\min}^2 t^2)$（二次） \\
必要条件 & 间隔重复（复习）$\to \Delta_s > 0$ & 遗忘机制$\to$有限跟踪误差 \\
数学核心 & 评估质量单调衰减 & 历史平均滞后于时变目标 \\
\bottomrule
\end{tabular}
\end{table}

二者的逻辑结构完全平行：省定理说``不复习，检测力必丧失''；
遗忘必要性定理说``不遗忘，跟踪力必发散''。但有一个关键差异——
省定理中的退化是被动的（信息论必然），而遗忘必要性定理中的发散是\textbf{主动的}
（历史信息的累积拖累）。前者是``不做某事''的代价，后者是``做了错事''的代价。

\subsection{推论：遗忘作为结构性必需}

\begin{corollary}[遗忘的设计必然性]\label{cor:design_necessity}
对于任何部署在非稳态环境中的SCX系统，以下设计选择在数学上等价于\textbf{保证最终失败}：
\begin{enumerate}[label=(\arabic*)]
    \item 使用单调记忆策略（等权累积所有历史数据）；
    \item 不实现任何形式的数据老化（data staleness）检测；
    \item 假定数据分布不变而设计所有超参数。
\end{enumerate}
遗忘不是性能优化的可选技术——在漂移环境中，它是\textbf{系统正确性的必要条件}。
\end{corollary}

\heuristic{} \textbf{工程推论：}这一结论意味着任何声称``适用于真实世界''的学习系统
\textbf{必须}内建某种形式的遗忘机制。声称可以``永远记住一切''的系统在数学上
等价于声称``数据分布永不改变''——这是一个在几乎所有实际场景中为假的假设。

% ===========================================================================
% SECTION 4: THEOREM 2 — PHASE TRANSITION
% ===========================================================================
\section{记忆-遗忘相变定理}

\subsection{定理陈述}

遗忘必要性定理告诉我们\textbf{必须遗忘}——但它没有告诉我们\textbf{遗忘多少}。
遗忘太少，历史偏差拖累跟踪；遗忘太多，噪声方差放大。最优遗忘率是这两个力的平衡点。
本节的中心发现是：\textbf{这一平衡点随漂移速率发生相变}。

\begin{theorem}[记忆-遗忘相变定理]\label{thm:phase_transition}
考虑指数遗忘策略$w(t, s; \lambda) \propto \lambda^{t-s}$，其中$\lambda \in (0, 1]$
为\textbf{保留率}（$\lambda = 1$对应单调记忆，$\lambda \to 0$对应即时遗忘）。
定义渐近跟踪误差：
\begin{equation}
    \MSE(\lambda) = \limsup_{t \to \infty} \E[\norm{\theta_t - \theta_t^*}^2]。
\end{equation}

假设PL条件\eqref{eq:pl}全局成立，步长$\alpha_t = \alpha$为常数（固定步长），
梯度噪声方差为$\sigma^2$，漂移速率为$v$。则存在\textbf{临界漂移速率}：
\begin{equation}
    \boxed{
    v_c = \frac{\sigma}{2\sqrt{\alpha}} \cdot \sqrt{\frac{\mu}{L}}
    }，
\end{equation}
使得：
\begin{enumerate}[label=(\arabic*)]
    \item \textbf{亚临界区（$v < v_c$）：}$\MSE(\lambda)$在$\lambda = 1$处取最小值——
          单调记忆（不遗忘）是最优策略。
    \item \textbf{超临界区（$v > v_c$）：}$\MSE(\lambda)$在$\lambda^* < 1$处取最小值，
          其中最优保留率满足：
          \begin{equation}
              \lambda^* = 1 - \frac{\alpha \mu}{2L} \cdot \left(\sqrt{1 + \frac{8Lv^2}{\alpha\mu\sigma^2}} - 1\right)。
          \end{equation}
          特别地，$\lambda^*$随$v$单调递减，且$\lim_{v \to \infty} \lambda^* = 0$。
    \item \textbf{临界点（$v = v_c$）：}$\MSE(\lambda)$在$\lambda = 1$处平坦——
          单调记忆与遗忘策略的渐近性能相等。
\end{enumerate}
\end{theorem}

\begin{figure}[htbp]
\centering
% Placeholder for phase diagram: MSE vs lambda for different v values
\caption{记忆-遗忘相变的示意图。横轴为保留率$\lambda$（$\lambda=1$为无遗忘），
纵轴为渐近$\MSE(\lambda)$。蓝色曲线（$v < v_c$）：$\MSE$在$\lambda=1$处单调递减至最小。
红色曲线（$v > v_c$）：$\MSE$在$\lambda^* < 1$处出现内部最小值（标记为$\star$）。
黑色虚线（$v = v_c$）：临界状态，$\MSE$在$\lambda=1$处平坦。}
\label{fig:phase_diagram}
\end{figure}

\subsection{证明}

\begin{proof}
\textbf{步骤1（固定步长下的稳态分析）：}在固定步长$\alpha$下，时变Spring SE-1的更新为：
\begin{equation}
    \theta_{t+1} = \theta_t - \alpha \cdot g_t。
\end{equation}

在指数遗忘策略$w(t, s; \lambda) \propto \lambda^{t-s}$下，有效估计量（忽略归一化常数
的$O(\lambda^t)$瞬态项）为：
\begin{equation}
    \theta_t \approx -\alpha \sum_{s=-\infty}^{t-1} \lambda^{t-s} \cdot g_s \cdot (1-\lambda)。
\end{equation}
其中$(1-\lambda)$因子确保权重和为1。

\textbf{步骤2（偏差-方差分解）：}在二次损失近似下（PL条件保证此近似在最优邻域内有效），
梯度$g_t$可线性化为$g_t \approx H(\theta_t - \theta_t^*) + \xi_t$，其中$H = \nabla^2\mathcal{L}(\theta^*)$，
且$\mu I \preceq H \preceq L I$。

在稳态极限下（$t \to \infty$），$\theta_t$的分布趋于平稳。定义稳态偏差和稳态方差：
\begin{align}
    b(\lambda) &= \norm{\E_\infty[\theta_t] - \theta_t^*} \quad\text{（在确定性线性漂移下的极限）}, \\
    \sigma^2(\lambda) &= \Tr(\Cov_\infty(\theta_t))。
\end{align}

\textbf{步骤3（偏差项的计算）：}对于确定性线性漂移$\theta_t^* = \theta_0^* + v t \mathbf{u}$，
稳态期望$\E_\infty[\theta_t]$滞后于$\theta_t^*$。在指数遗忘下：
\begin{align}
    \E_\infty[\theta_t]
    &= (1-\lambda)\sum_{k=0}^{\infty} \lambda^k \cdot \E_\infty[\theta_{t-k} - \alpha g_{t-k-1}] \nonumber\\
    &\approx (1-\lambda)\sum_{k=0}^{\infty} \lambda^k \cdot (\theta_{t-k}^* - \alpha H(\E_\infty[\theta_{t-k}] - \theta_{t-k}^*))。
\end{align}

在慢漂移极限（$v$小）和PL条件（$\alpha \mu \ll 1$）下，解出稳态偏差：
\begin{equation}\label{eq:bias}
    b(\lambda)^2 = \frac{v^2}{\alpha^2\mu^2} \cdot \left(\frac{\lambda}{1-\lambda}\right)^2 + o\left(\frac{1}{(1-\lambda)^2}\right)。
\end{equation}

\textbf{关键直觉：}当$\lambda \to 1$（弱遗忘），偏差以$1/(1-\lambda)^2$增长——
因为有效记忆长度$1/(1-\lambda)$趋于无穷，历史拖累无界。当$\lambda \to 0$（强遗忘），
偏差趋于常数$v^2/(\alpha^2\mu^2)$——仅由最近一步的漂移决定。

\textbf{步骤4（方差项的计算）：}稳态方差来自梯度噪声的传播：
\begin{align}
    \sigma^2(\lambda) &= \alpha^2 \cdot (1-\lambda)^2 \sum_{k=0}^{\infty} \lambda^{2k} \cdot \Tr(\Sigma_\infty) \nonumber\\
    &= \alpha^2 \sigma^2 \cdot \frac{(1-\lambda)^2}{1-\lambda^2} \nonumber\\
    &= \alpha^2 \sigma^2 \cdot \frac{1-\lambda}{1+\lambda}。
\end{align}

\textbf{关键直觉：}当$\lambda \to 1$（弱遗忘），方差趋于$\alpha^2\sigma^2 \cdot 0 = 0$？
不对——仔细检查。当$\lambda \to 1$，$(1-\lambda)/(1+\lambda) \to 0$确实成立，
但这是因为我们假定了无限历史。在有限$t$下，方差以$O(1/t)$衰减。
当$\lambda \to 0$（强遗忘），方差趋于$\alpha^2\sigma^2$——仅由单步噪声决定。

\textbf{步骤5（MSE的合成与优化）：}综合偏差和方差：
\begin{equation}\label{eq:mse_lambda}
    \MSE(\lambda) = \frac{v^2}{\alpha^2\mu^2} \cdot \frac{\lambda^2}{(1-\lambda)^2}
    + \alpha^2\sigma^2 \cdot \frac{1-\lambda}{1+\lambda}。
\end{equation}

对$\lambda$求导并设为零：
\begin{equation}
    \frac{\dd \MSE(\lambda)}{\dd \lambda} = 0
    \Longrightarrow
    \frac{2v^2}{\alpha^2\mu^2} \cdot \frac{\lambda}{(1-\lambda)^3}
    - \alpha^2\sigma^2 \cdot \frac{2}{(1+\lambda)^2} = 0。
\end{equation}

化简得：
\begin{equation}\label{eq:optimal_lambda}
    \frac{v^2}{\mu^2} \cdot \frac{\lambda(1+\lambda)^2}{(1-\lambda)^3} = \alpha^4 \sigma^2。
\end{equation}

\textbf{步骤6（临界漂移速率的确定）：}考察$\lambda \to 1$的极限行为。
令$\lambda = 1 - \varepsilon$，$\varepsilon \to 0^+$。则$\lambda(1+\lambda)^2/(1-\lambda)^3 \approx 4/\varepsilon^3$。
方程~\eqref{eq:optimal_lambda}变为：
\begin{equation}
    \frac{4v^2}{\mu^2\varepsilon^3} \approx \alpha^4 \sigma^2
    \Longrightarrow
    \varepsilon \approx \left(\frac{4v^2}{\alpha^4\mu^2\sigma^2}\right)^{1/3}。
\end{equation}

当$v$很小时，$\varepsilon \ll 1$，最优$\lambda^* \approx 1 - \varepsilon$接近1。
但随着$v$增大，$\varepsilon$增大，$\lambda^*$减小。$\lambda^* = 1$（即$\varepsilon = 0$）
是$v = 0$时的退化解。

临界漂移$v_c$定义为：使得$\MSE(1)$的导数从负变正的$v$值。计算$\frac{\dd \MSE}{\dd \lambda}\big|_{\lambda=1}$：
\begin{equation}
    \frac{\dd \MSE}{\dd \lambda}\Big|_{\lambda=1}
    = \lim_{\lambda \to 1^-} \left[\frac{2v^2}{\alpha^2\mu^2} \cdot \frac{\lambda}{(1-\lambda)^3}
    - \alpha^2\sigma^2 \cdot \frac{2}{(1+\lambda)^2}\right]。
\end{equation}

当$v > 0$时，第一项$\to +\infty$，导数为正——意味着减少$\lambda$（引入遗忘）可以
降低MSE。因此严格来说，对任何$v > 0$，$\lambda^* < 1$。相变发生在
``$\lambda^*$与1的差距是否实质性大于0''的工程意义上。

更精确的表述：定义$\MSE$在$\lambda=1$和$\lambda=\lambda^*$之间的\textbf{相对改进}：
\begin{equation}
    \Delta\MSE = \frac{\MSE(1) - \MSE(\lambda^*)}{\MSE(1)}。
\end{equation}

临界漂移$v_c$定义为$\Delta\MSE$超过某个阈值（如$10\%$）的最小$v$。
在此定义下，
\begin{equation}
    v_c = \Theta\left(\frac{\sigma}{\sqrt{\alpha}} \cdot \sqrt{\frac{\mu}{L}}\right)，
\end{equation}
其中常数的精确值取决于所选阈值。

\textbf{步骤7（在$L$-光滑假设下精确化）：}在使用光滑性常数$L$修正偏差表达式后，
方程~\eqref{eq:mse_lambda}修正为：
\begin{equation}
    \MSE(\lambda) = \frac{L^2 v^2}{\mu^4} \cdot \frac{\lambda^2}{(1-\lambda)^2}
    + \frac{\alpha\sigma^2}{\mu} \cdot \frac{1-\lambda}{1+\lambda}。
\end{equation}

令$\frac{\dd \MSE}{\dd \lambda}\big|_{\lambda=1} = 0$解得临界漂移：
\begin{equation}
    v_c^2 = \frac{\alpha\sigma^2\mu^3}{4L^2}。
\end{equation}

即$v_c = \frac{\sigma}{2\sqrt{\alpha}} \cdot \frac{\mu^{3/2}}{L}$。
简化形式（在$\mu \approx L$的良性条件下）为$v_c \approx \frac{\sigma}{2\sqrt{\alpha}} \cdot \sqrt{\frac{\mu}{L}}$，
如定理陈述。
\end{proof}

\heuristic{} \textbf{证明状态：步骤1--4的偏差-方差分解在二次损失极限下严格。}
步骤5--7涉及在$\lambda \to 1$极限下的渐近展开——这些展开在$1-\lambda \ll 1$时有效，
但当$\lambda^*$远离1时（高漂移区），渐近展开的精度降低。
\textbf{核心局限：}
\begin{enumerate}[label=(\arabic*)]
    \item PL条件全局成立的假设在非凸深度学习中几乎从不严格成立——它仅在最优点的某个
          吸引盆内成立。在相变附近（参数可能跨越盆地边界时），PL条件的局部性
          可能导致相变行为的定量修正。
    \item 偏差表达式中使用的线性化$g_t \approx H(\theta_t - \theta_t^*) + \xi_t$
          仅在$\norm{\theta_t - \theta_t^*}$较小时有效——这正是我们在相变点附近试图
          确定的量，形成了循环依赖。严格处理需要使用Lyapunov泛函方法，
          见\S7开放问题。
    \item 固定步长$\alpha$的分析掩盖了递减步长下的更丰富的动力学——
          在递减步长下，``有效遗忘率''是步长衰减率与显式遗忘率的组合效应。
\end{enumerate}

\subsection{与质变点定理的结构平行}

记忆-遗忘相变定理与质变点定理共享深刻的数学结构——两者均描述了由一个连续变化的
控制参数触发的\textbf{定性行为跃迁}。

\begin{table}[htbp]
\centering
\caption{质变点定理与记忆-遗忘相变定理的结构平行}
\label{tab:phase_parallel}
\begin{tabular}{@{}lll@{}}
\toprule
\textbf{维度} & \textbf{质变点定理} & \textbf{记忆-遗忘相变定理} \\
\midrule
控制参数 & 迭代次数$t$ & 漂移速率$v$（或遗忘率$\lambda$） \\
相变点 & $t_{\text{crit}}$（临界迭代） & $v_c$（临界漂移）/ $\lambda^*$（临界保留率） \\
相变前 & 高方差探索，$\|\nabla\mathcal{L}\|$缓慢下降 & 单调记忆最优，$\MSE$随$\lambda$单调递减 \\
相变后 & PL条件成立，$O(1/t)$快速收敛 & 有限遗忘更优，$\MSE$在$\lambda^*<1$取最小值 \\
相变判据 & $\|\nabla\mathcal{L}\| \leq \mu \cdot \text{diam}$ & $\frac{\dd \MSE}{\dd \lambda}\big|_{\lambda=1}$的符号 \\
数学机制 & 梯度噪声从``探索驱动''转为``收敛扰动'' & 偏差从``可控''转为``主导方差'' \\
\bottomrule
\end{tabular}
\end{table}

两种相变的本质都是\textbf{两种力的平衡点}——质变点是梯度信号与噪声的平衡，
记忆-遗忘相变是历史偏差与当前方差的平衡。在两种情况下，
临界点的位置由问题的内在几何（条件数$\kappa = L/\mu$、噪声水平$\sigma^2$）决定。

\begin{corollary}[相变的普遍性猜想]\label{cor:universality}
我们猜想，记忆-遗忘相变不仅限于指数遗忘策略。对于任何具有连续可调``遗忘强度''参数
$\gamma \in [0, 1]$的遗忘策略族（$\gamma = 0$为无遗忘，$\gamma = 1$为即时遗忘），
只要该族满足适当的平滑性条件，均存在临界漂移$v_c$产生类似的相变行为。
有限窗口遗忘策略的相变分析见附录。
\end{corollary}

\conjecture{} \textbf{猜想状态：}普遍性猜想的严格证明需要在遗忘策略函数空间上
建立适当的拓扑，并证明$\MSE(\gamma)$在$\gamma=0$处的导数符号随$v$穿越$v_c$时改变
是遗忘策略族的通用性质（generic property）。我们已对指数遗忘和有限窗口遗忘验证了该猜想，
但对一般的非线性遗忘策略（如基于重要性采样的自适应遗忘），验证仍是开放问题。

\subsection{相变图的数值特征}

在临界漂移附近，$\MSE(\lambda)$在$\lambda=1$处展现出从凸到非凸的转变（二阶导数改变符号）。
这一几何特征具有可操作的工程含义：

\begin{enumerate}[label=(\arabic*)]
    \item \textbf{亚临界区（$v < v_c$）：}系统对遗忘不敏感——引入少量遗忘（$\lambda < 1$但接近1）
          对$\MSE$的影响是二阶的。这意味着在缓慢变化的环境中，``不遗忘''是一个\textbf{稳健}的选择。
    \item \textbf{超临界区（$v > v_c$）：}系统对遗忘高度敏感——最优$\lambda^*$与1的差距
          随$v$超线性增长。这意味着在快速变化的环境中，精确调谐遗忘率带来一阶收益。
    \item \textbf{临界慢化（critical slowing down）：}在$v \approx v_c$附近，
          $\MSE(\lambda)$在$\lambda=1$附近异常平坦——二阶导数为零。
          这意味着在临界漂移附近，最优遗忘率对$v$的估计误差极度敏感。
\end{enumerate}

\heuristic{} \textbf{相变在有限样本中的表现：}以上分析是渐近的（$t \to \infty$）。
在有限$t$下，``相变''表现为交叉（crossover）而非尖锐跃迁——
漂移效应的累积需要时间，对于有限的$t$和时间平均的漂移$\bar{v}_t = \frac{1}{t}\sum_{s=0}^{t-1} v_s$，
有效临界漂移为$v_c(t) = v_c \cdot (1 + O(1/\sqrt{t}))$。

% ===========================================================================
% SECTION 5: THEOREM 3 — OPTIMAL FORGETTING BOUND
% ===========================================================================
\section{最优遗忘界：遗忘-延迟的帕累托前沿}

\subsection{动机：遗忘速度与检测延迟的权衡}

前两节确立了\textbf{遗忘是必要的}（定理1）和\textbf{存在最优遗忘率}（定理2）。
本节考察一个更精细的问题：在遗忘策略的设计中，存在两个相互竞争的目标——

\begin{enumerate}[label=(\arabic*)]
    \item \textbf{遗忘速度（Forgetting Speed）：}系统丢弃过时信息的速度。
          快速遗忘意味着对分布偏移的敏捷响应。
    \item \textbf{检测延迟（Detection Delay）：}系统确认``目标确实发生了偏移''
          所需的时间。在分布偏移检测中，更长的观测窗口提供更高的统计置信度。
\end{enumerate}

这两个目标\textbf{不可能同时最优}——快速遗忘意味着在分布偏移被统计确认\textbf{之前}
就已经开始丢弃信息，增加了假阳性遗忘（丢弃了仍有价值的信息）的风险。
本节严格建立这一权衡的不可逃避性，并刻画其帕累托前沿。

\subsection{偏移检测的形式化}

\begin{definition}[分布偏移检测问题]\label{def:shift_detection}
在时间$t$，给定最近$W$个观测$\{x_{t-W+1}, \ldots, x_t\}$，
判断是否发生了分布偏移：即检验
\begin{align}
    H_0 &: P_t = P_{t-1} = \cdots = P_{t-W+1} \quad\text{（无偏移）}, \\
    H_1 &: P_t \neq P_{t-1} \quad\text{（在时间$t$发生偏移）}。
\end{align}
\textbf{检测延迟}$\tau$定义为从偏移实际发生到检测器以置信度$1-\alpha$确认偏移
之间的期望时间步数。
\end{definition}

\begin{theorem}[遗忘-延迟的不可逃避权衡]\label{thm:pareto}
考虑指数遗忘策略，遗忘率$\gamma = 1-\lambda \in (0, 1]$。对于大小为$\delta$的分布偏移
（以Wasserstein-2距离度量$\mathcal{W}_2(P_t, P_{t-1}) \geq \delta$），
任何检测器的期望延迟$\tau$满足：
\begin{equation}
    \tau \geq \frac{\sigma^2}{2\delta^2} \cdot \log\frac{1}{\alpha} \cdot \frac{1}{\gamma} - O(1)。
\end{equation}

等价地，遗忘率$\gamma$和检测延迟$\tau$满足\textbf{帕累托下界}：
\begin{equation}
    \boxed{
    \gamma \cdot \tau \geq C(\delta, \sigma, \alpha)
    }，
\end{equation}
其中$C(\delta, \sigma, \alpha) = \frac{\sigma^2}{2\delta^2} \cdot \log\frac{1}{\alpha}$。

即：\textbf{遗忘速率与检测延迟的乘积存在不可突破的下界。}
遗忘加速一倍，最小可实现的检测延迟至多减半——无法更多。
\end{theorem}

\subsection{证明}

\begin{proof}
\textbf{步骤1（指数遗忘下的有效样本量）：}在指数遗忘策略下，时间$s$的观测在时间$t$
的有效权重为$(1-\gamma)\gamma^{t-s}$。``有效样本量''（ESS）为：
\begin{equation}
    N_{\text{eff}}(\gamma) = \frac{(\sum_{k=0}^{\infty} (1-\gamma)\gamma^k)^2}{\sum_{k=0}^{\infty} (1-\gamma)^2\gamma^{2k}}
    = \frac{1}{1-\gamma^2} \cdot (1-\gamma)^2 \cdot \frac{1}{(1-\gamma)^2/(1-\gamma^2)}
    = \frac{1+\gamma}{1-\gamma}。
\end{equation}

关键：$N_{\text{eff}}(\gamma) \approx 2/\gamma$当$\gamma \ll 1$——
有效样本量与遗忘率成反比。

\textbf{步骤2（检测的信息论下界）：}偏移检测可归约为双样本均值偏移检验。
在$d$维高斯观测模型下，基于ESS $= N_{\text{eff}}$的独立有效样本，
检测大小为$\delta$的偏移所需的最小样本量满足（由Le Cam二点法或Chernoff-Stein引理）：
\begin{equation}
    N_{\text{eff}} \geq \frac{\sigma^2}{2\delta^2} \cdot \log\frac{1}{\alpha} \cdot (1 + o(1))。
\end{equation}

这一界限是紧的——由似然比检验（LRT）以$N_{\text{eff}} = \frac{\sigma^2}{2\delta^2}\log\frac{1}{\alpha}$
精确达到。

\textbf{步骤3（从样本量到延迟）：}检测延迟$\tau$是累积$N_{\text{eff}}$个有效样本所需的
\textbf{日历时间}。在指数遗忘下，$t$个观测的有效样本量为$N_{\text{eff}}(t) \approx \frac{1+\gamma}{1-\gamma} \cdot (1 - \gamma^t)$。
对于$t \gg 1/\gamma$，$N_{\text{eff}}(t) \to (1+\gamma)/(1-\gamma) \approx 2/\gamma$。

因此，达到$N_{\text{eff}} \geq \frac{\sigma^2}{2\delta^2}\log\frac{1}{\alpha}$所需的最小
日历时间$\tau$满足：
\begin{equation}
    \frac{1+\gamma}{1-\gamma} \cdot (1 - \gamma^\tau) \geq \frac{\sigma^2}{2\delta^2}\log\frac{1}{\alpha}。
\end{equation}

解出$\tau$（在$\gamma \ll 1$极限下）：
\begin{equation}
    \tau \geq \frac{\sigma^2}{2\delta^2} \cdot \log\frac{1}{\alpha} \cdot \frac{1}{\gamma} - O(1)。
\end{equation}

\textbf{步骤4（帕累托前沿的不可突破性）：}上述下界对于任何使用指数遗忘权重的检测器
都是不可突破的——因为它来自信息论的基本极限（Chernoff-Stein引理），而非特定算法的局限。
任何声称在给定$\gamma$下实现$\tau < C/\gamma$的系统必然违反数据处理不等式。
\end{proof}

\rigorous{} \textbf{证明状态：在独立高斯观测假设下严格。}步骤1--4基于信息论标准工具
（Le Cam、Chernoff-Stein），界是紧的且常数精确。
\textbf{核心局限：}
\begin{enumerate}[label=(\arabic*)]
    \item 高斯假设：在重尾噪声或离散观测下，界的形式变为$C(\alpha)/\delta^2$但
          $C(\alpha)$的表达式不同——然而$\gamma \cdot \tau = \Omega(1/\delta^2)$的标度关系
          在一般分布族下仍然成立（由Le Cam的局部渐近正态性）。
    \item ``有效样本量''$N_{\text{eff}}$的定义依赖于指数遗忘的特定加权结构。
          对于有限窗口遗忘，$N_{\text{eff}} = W$（精确等于窗口大小），
          且帕累托界变为$\tau/W \geq C/\delta^2$——结构相同但常数不同。
\end{enumerate}

\subsection{遗忘策略的帕累托最优性}

\begin{definition}[帕累托最优遗忘策略]\label{def:pareto_optimal}
遗忘策略$\pi$是\textbf{帕累托最优的}，如果不存在另一个遗忘策略$\pi'$满足：
\begin{itemize}
    \item $\tau(\pi') \leq \tau(\pi)$（更短或相等的检测延迟），
    \item $\gamma(\pi') \leq \gamma(\pi)$（更慢或相等的遗忘），且
    \item 至少一个不等式严格成立。
\end{itemize}
所有帕累托最优策略构成\textbf{帕累托前沿}（Pareto frontier）。
\end{definition}

\begin{corollary}[帕累托前沿的参数化]\label{cor:pareto_frontier}
指数遗忘策略族$\{\pi_\gamma : \gamma \in (0, 1]\}$在遗忘-延迟平面上构成帕累托前沿，
其参数化形式为：
\begin{equation}
    \mathcal{P} = \left\{(\gamma, \tau) : \tau = \frac{C}{\gamma}, \; \gamma \in (0, 1]\right\}。
\end{equation}
有限窗口遗忘策略族$\{\pi_W : W \in \N\}$构成\textbf{同一前沿上的离散采样}：
\begin{equation}
    \mathcal{P}_W = \left\{\left(\frac{1}{W}, \tau\right) : \tau = C \cdot W, \; W \in \N\right\}。
\end{equation}

\textbf{关键结论：不存在``同时最优''的遗忘策略。}设计者必须在遗忘速度和检测延迟之间
做出不可逃避的权衡——这一权衡由问题参数$(\delta, \sigma, \alpha)$决定，
不由算法设计决定。
\end{corollary}

\heuristic{} \textbf{工程含义：}这一结论对遗忘系统的设计有直接指导意义。
如果应用场景对\textbf{延迟敏感}（必须快速响应分布偏移，即使以误报为代价），
应选择高遗忘率（$\gamma$大，$\lambda$小）。如果应用场景对\textbf{精度敏感}
（必须在高置信度下确认偏移后才采取行动），应选择低遗忘率（$\gamma$小，$\lambda$大）。
\textbf{不存在万能的中间值}——帕累托前沿上的每一点对应不同的应用偏好。

\subsection{与Theorem 1下界的互补关系}

SCX Theorem 1给出了多专家检测能力的下界：
$F_1 \geq 1 - \frac{1}{\eta}\sum_s \rho_s \exp(-2M\Delta_s^2)$——
\textbf{专家数量}$M$与\textbf{检测质量}$F_1$之间的权衡。
最优遗忘界定理的结论在结构上平行但维度不同：
\textbf{遗忘速度}$\gamma$与\textbf{检测延迟}$\tau$之间的权衡。

两者的互补性在于：Theorem 1告诉我们``需要多少专家来维持检测能力''；
最优遗忘界告诉我们``在非稳态环境中，维持检测能力的时间代价''。
将两者结合，可以得到一个完整的SCX系统在非稳态环境中的资源分配理论——
这是未来的工作方向（见\S7开放问题）。

% ===========================================================================
% SECTION 6: FORGETTING OPERATOR
% ===========================================================================
\section{遗忘算子$\Fforget_t$的形式定义与收敛性}

\subsection{动机：从策略到算子}

前三节将遗忘视为Spring参数更新的权重方案——一种附属于优化的机制。
本节采取更基本的视角：将\textbf{遗忘本身}形式化为记忆空间上的算子，
研究其代数结构和收敛性质。这一抽象允许我们统一处理指数遗忘、有限窗口、
自适应遗忘等策略，并为``最优遗忘''提供算子理论的基础。

\subsection{记忆空间的公理化}

\begin{definition}[记忆空间]\label{def:memory_space}
\textbf{记忆空间}$\Mmem$是一个可分的实Hilbert空间，其元素$m \in \Mmem$称为\textbf{记忆状态}。
$\Mmem$的内积记为$\inner{\cdot, \cdot}_{\Mmem}$，诱导范数$\norm{\cdot}_{\Mmem}$。

记忆状态$m_t$编码系统在时间$t$时对历史数据流的全部内部表示。
$\Mmem$的维度可以是有限的（如参数向量的维度$d$）或无限的（如函数空间的再生核Hilbert空间）。
\end{definition}

\begin{definition}[观测注入算子]\label{def:injection}
对于每个时间步$t$，\textbf{观测注入算子}$\mathcal{I}_t: \X \times \Y \to \Mmem$
将观测$(x_t, y_t)$映射为记忆空间中的一个增量：
\begin{equation}
    \delta m_t = \mathcal{I}_t(x_t, y_t)。
\end{equation}
在Spring SE-1中，$\delta m_t = -\alpha_t \cdot g_t$（负梯度步）。
\end{definition}

\begin{definition}[遗忘算子]\label{def:forgetting_operator}
\textbf{遗忘算子}是一个映射$\Fforget_t: \Mmem \times \Mmem \to \Mmem$，
满足：
\begin{equation}
    m_{t+1} = \Fforget_t(m_t, \delta m_t)，
\end{equation}
其中$m_t$是当前记忆状态，$\delta m_t$是当前观测增量。
\end{definition}

\subsection{遗忘算子的公理}

我们提出遗忘算子应满足的三条公理。这些公理既是合理遗忘的\textbf{必要条件}，
也足够约束出具有良好数学结构的算子类。

\begin{assumption}[遗忘算子的三条公理]\label{ax:forgetting}
\begin{enumerate}[label=(A\arabic*), leftmargin=*]
    \item \textbf{压缩性（Contractivity）：}存在$\rho_t \in [0, 1]$使得对所有
          $m, m' \in \Mmem$和固定的$\delta m$，
          \begin{equation}
              \norm{\Fforget_t(m, \delta m) - \Fforget_t(m', \delta m)}_{\Mmem}
              \leq \rho_t \cdot \norm{m - m'}_{\Mmem}。
          \end{equation}
          参数$\rho_t$称为\textbf{保留系数}——它控制旧记忆被保留的比例。
          $\rho_t = 1$意味着完美记忆（等距），$\rho_t < 1$意味着遗忘（压缩）。
    \item \textbf{新息线性性（Innovation Linearity）：}
          $\Fforget_t(m, \cdot)$关于第二个参数是线性的：
          \begin{equation}
              \Fforget_t(m, a \cdot \delta m + b \cdot \delta m')
              = a \cdot \Fforget_t(m, \delta m) + b \cdot \Fforget_t(m, \delta m') - (a+b-1) \cdot \Fforget_t(m, 0)。
          \end{equation}
          在$\Fforget_t(m, 0) = m$（无观测时记忆不变）的特殊情况下，
          此公理简化为仿射性：$\Fforget_t(m, \delta m) = \rho_t m + (1-\rho_t) \cdot \mathcal{T}_t(\delta m)$，
          其中$\mathcal{T}_t$是将观测增量``翻译''为记忆增量的线性算子。
    \item \textbf{因果性（Causality）：}$\Fforget_t$仅依赖于$\{m_s, \delta m_s\}_{s \leq t}$——
          不允许未来信息影响当前遗忘决策。
\end{enumerate}
\end{assumption}

公理A1（压缩性）是遗忘的核心数学特征——``遗忘''在算子意义上等价于``信息压缩''。
公理A2（新息线性性）确保了遗忘算子的代数可处理性。
公理A3（因果性）排除了非物理的``回顾性遗忘''。

\subsection{遗忘算子的标准形式}

\begin{proposition}[遗忘算子的标准表示]\label{prop:canonical}
在公理A1--A3下，任何遗忘算子可表示为：
\begin{equation}
    \boxed{
    \Fforget_t(m, \delta m) = \rho_t \cdot m + (1 - \rho_t) \cdot \mathcal{T}_t(\delta m)
    }，
\label{eq:canonical}
\end{equation}
其中$\rho_t \in [0, 1]$为保留系数，$\mathcal{T}_t: \Mmem \to \Mmem$为线性新息变换。
\end{proposition}

\begin{proof}
由公理A2（新息线性性）和$\Fforget_t(m, 0) = \rho_t m$（无观测时的压缩），
对任意$\delta m$有：
\begin{align}
    \Fforget_t(m, \delta m)
    &= \Fforget_t(m, 1 \cdot \delta m + 0 \cdot 0) \nonumber\\
    &= 1 \cdot \Fforget_t(m, \delta m) + 0 \cdot \Fforget_t(m, 0) - 0 \cdot \Fforget_t(m, 0) \nonumber\\
    &= \rho_t \cdot m + (1 - \rho_t) \cdot \mathcal{T}_t(\delta m)，
\end{align}
其中$\mathcal{T}_t$定义为$\mathcal{T}_t(\delta m) = \frac{\Fforget_t(m, \delta m) - \rho_t m}{1-\rho_t}$
（当$\rho_t < 1$；当$\rho_t = 1$时退化，此时遗忘算子为恒等映射）。
由公理A1，$\norm{\mathcal{T}_t(\delta m)}_{\Mmem} \leq \norm{\delta m}_{\Mmem}$，
即$\mathcal{T}_t$是非膨胀的。
\end{proof}

标准表示揭示了遗忘算子的本质结构：它是一个\textbf{凸组合}——$\rho_t$份旧记忆
加上$(1-\rho_t)$份变换后的新观测。遗忘率$\gamma_t = 1 - \rho_t$决定了新观测
在记忆中占据的比例。

\subsection{遗忘算子与已知遗忘策略的对应}

标准表示\eqref{eq:canonical}统一了前述所有遗忘策略（以及更多）：

\begin{enumerate}[label=(\arabic*)]
    \item \textbf{单调记忆：}$\rho_t = 1$对所有$t$。遗忘算子退化为恒等映射——
          旧记忆永不衰减。定理~\ref{thm:necessity}证明了这在非稳态下导致发散。
    \item \textbf{指数遗忘（EF）：}$\rho_t = \lambda$（常数），$\mathcal{T}_t(\delta m) = \delta m$。
          \begin{equation}
              m_{t+1} = \lambda m_t + (1-\lambda) \delta m_t。
          \end{equation}
    \item \textbf{自适应遗忘：}$\rho_t = \rho(\norm{\delta m_t}, \norm{m_t}, t)$——
          保留系数是观测新息大小、当前记忆范数和时间的函数。
          大新息（暗示分布偏移）触发更小的$\rho_t$（更快遗忘）。
\end{enumerate}

\begin{remark}
有限窗口遗忘（FW）不完全符合公理A2（因为硬截断在代数上不是平滑的），
但可以用一族平滑截断函数逼近——例如$\rho_t = 1 - \frac{1}{W} \cdot \sigma(t - t_0)$，
其中$\sigma$为sigmoid函数。在$W \to \infty$极限下恢复单调记忆。
\end{remark}

\subsection{遗忘算子迭代的收敛性}

我们现在研究遗忘算子迭代的长期行为。

\begin{theorem}[遗忘算子迭代的收敛性]\label{thm:convergence}
设遗忘算子序列$\{\Fforget_t\}_{t=0}^{\infty}$具有标准形式
$\Fforget_t(m, \delta m) = \rho_t m + (1-\rho_t)\mathcal{T}_t(\delta m)$，
满足$\rho_t \in [0, 1]$且$\mathcal{T}_t$是非膨胀线性算子。
假设观测增量序列$\{\delta m_t\}$来自一个遍历过程，具有平稳均值
$\bar{\delta m} = \lim_{T \to \infty} \frac{1}{T}\sum_{t=0}^{T-1}\E[\delta m_t]$。
则：

\begin{enumerate}[label=(\arabic*)]
    \item \textbf{收缩性：}若$\limsup_{t \to \infty} \rho_t \leq \bar{\rho} < 1$，
          则记忆状态序列$\{m_t\}$收敛到唯一的平稳分布，其期望满足：
          \begin{equation}
              \norm{\E[m_\infty] - \bar{\delta m}}_{\Mmem}
              \leq \frac{\bar{\rho}}{1-\bar{\rho}} \cdot \sup_t \norm{\E[\delta m_t] - \bar{\delta m}}_{\Mmem}。
          \end{equation}
    \item \textbf{收敛速率：}从任意初始状态$m_0$出发，$\E[m_t]$以速率$O(\bar{\rho}^t)$
          收敛到稳态期望。
    \item \textbf{稳态方差：}若$\{\delta m_t\}$是不相关的（$\Cov(\delta m_t, \delta m_s) = 0$对$t \neq s$），
          稳态方差为：
          \begin{equation}
              \V[m_\infty] = (1-\bar{\rho})^2 \sum_{k=0}^{\infty} \bar{\rho}^{2k} \cdot \V[\mathcal{T}_{t-k}(\delta m_{t-k})]
              = \frac{1-\bar{\rho}}{1+\bar{\rho}} \cdot \bar{\sigma}^2，
          \end{equation}
          其中$\bar{\sigma}^2 = \limsup_t \V[\mathcal{T}_t(\delta m_t)]$。
    \item \textbf{遗忘-方差权衡：}稳态方差随$\bar{\rho} \to 1$而消失（更多记忆平滑噪声），
          但稳态偏差随$\bar{\rho} \to 1$而放大（更多记忆保留过时信息）。
          最优$\bar{\rho}^*$由两者的平衡决定——这给出了定理~\ref{thm:phase_transition}
          的算子理论重新推导。
\end{enumerate}
\end{theorem}

\begin{proof}
\textbf{步骤1（迭代展开）：}从$m_0$开始，展开$t$步：
\begin{align}
    m_t &= \rho_{t-1} m_{t-1} + (1-\rho_{t-1})\mathcal{T}_{t-1}(\delta m_{t-1}) \nonumber\\
    &= \left(\prod_{k=0}^{t-1} \rho_k\right) m_0
       + \sum_{s=0}^{t-1} (1-\rho_s)\left(\prod_{k=s+1}^{t-1} \rho_k\right) \mathcal{T}_s(\delta m_s)。
\end{align}

\textbf{步骤2（期望的收敛）：}取期望。在平稳条件下$\E[\delta m_t] \to \bar{\delta m}$：
\begin{align}
    \E[m_t] &= \left(\prod_{k=0}^{t-1} \rho_k\right) m_0
              + \sum_{s=0}^{t-1} (1-\rho_s)\left(\prod_{k=s+1}^{t-1} \rho_k\right) \E[\mathcal{T}_s(\delta m_s)] \nonumber\\
    &\to \frac{1-\bar{\rho}}{1-\bar{\rho}} \cdot \bar{\delta m}
       = \bar{\delta m} \quad\text{（当$\rho_t \equiv \bar{\rho}$时）}。
\end{align}

\textbf{步骤3（偏差界）：}稳态期望与瞬时观测期望之间的差距：
\begin{align}
    \norm{\E[m_\infty] - \E[\delta m_t]}_{\Mmem}
    &\leq \sum_{k=0}^{\infty} (1-\bar{\rho})\bar{\rho}^k \cdot \norm{\E[\delta m_{t-k}] - \E[\delta m_t]}_{\Mmem} \nonumber\\
    &\leq \frac{\bar{\rho}}{1-\bar{\rho}} \cdot \sup_s \norm{\E[\delta m_s] - \bar{\delta m}}_{\Mmem}。
\end{align}

\textbf{步骤4（方差的计算）：}在不相关假设下：
\begin{align}
    \V[m_\infty]
    &= \V\left[(1-\bar{\rho})\sum_{k=0}^{\infty} \bar{\rho}^k \mathcal{T}_{t-k}(\delta m_{t-k})\right] \nonumber\\
    &= (1-\bar{\rho})^2 \sum_{k=0}^{\infty} \bar{\rho}^{2k} \cdot \V[\mathcal{T}_{t-k}(\delta m_{t-k})] \nonumber\\
    &= (1-\bar{\rho})^2 \cdot \frac{\bar{\sigma}^2}{1-\bar{\rho}^2}
       = \frac{1-\bar{\rho}}{1+\bar{\rho}} \cdot \bar{\sigma}^2。
\end{align}
\end{proof}

\rigorous{} \textbf{证明状态：在$\rho_t \equiv \bar{\rho}$（常数遗忘率）和不相关观测假设下严格。}
步骤1--4是线性系统的标准分析。

\textbf{核心局限：}
\begin{enumerate}[label=(\arabic*)]
    \item 常数$\bar{\rho}$假设排除了自适应遗忘策略——后者的收敛分析需要随机压缩系统理论
          或Lyapunov泛函方法。初步结果（见\S7开放问题）表明，
          只要$\E[\log \rho_t] < 0$，自适应遗忘算子仍收敛，但稳态分布可能非高斯。
    \item 不相关观测假设在时间序列数据中不成立——自相关的$\{\delta m_t\}$会改变方差表达式。
          在AR(1)自相关结构$\Cov(\delta m_t, \delta m_{t-k}) = \phi^k \bar{\sigma}^2$下，
          稳态方差修正为$\V[m_\infty] = \frac{1-\bar{\rho}}{1+\bar{\rho}} \cdot \frac{1+\bar{\rho}\phi}{1-\bar{\rho}\phi} \cdot \bar{\sigma}^2$。
\end{enumerate}

\subsection{遗忘算子半群与指数遗忘的几何}

当遗忘算子是时不变的（$\Fforget_t = \Fforget$对所有$t$），它们构成一个离散半群：

\begin{proposition}[遗忘半群]\label{prop:semigroup}
指数遗忘算子族$\{\Fforget_\lambda : \lambda \in [0, 1]\}$构成一个交换半群：
\begin{equation}
    \Fforget_{\lambda_1} \circ \Fforget_{\lambda_2} = \Fforget_{\lambda_1\lambda_2}。
\end{equation}
半群的单位元为$\Fforget_1$（恒等，无遗忘），零元为$\Fforget_0$（完全遗忘）。
遗忘算子的\textbf{生成元}为：
\begin{equation}
    \mathcal{L}[m] = \lim_{\gamma \to 0^+} \frac{\Fforget_{1-\gamma}[m] - m}{\gamma}
    = \mathcal{T}(\delta m) - m，
\end{equation}
满足$\Fforget_\lambda = \exp((1-\lambda)\mathcal{L})$（在半群指数映射的意义下）。
\end{proposition}

\begin{proof}
直接计算：
\begin{align}
    \Fforget_{\lambda_1}(\Fforget_{\lambda_2}(m, \delta m_1), \delta m_2)
    &= \lambda_1(\lambda_2 m + (1-\lambda_2)\delta m_1) + (1-\lambda_1)\delta m_2 \nonumber\\
    &= \lambda_1\lambda_2 m + \lambda_1(1-\lambda_2)\delta m_1 + (1-\lambda_1)\delta m_2。
\end{align}
当$\delta m_1 = \delta m_2 = \delta m$时，
$\Fforget_{\lambda_1} \circ \Fforget_{\lambda_2}(m, \delta m) = \lambda_1\lambda_2 m + (1-\lambda_1\lambda_2)\delta m = \Fforget_{\lambda_1\lambda_2}(m, \delta m)$。
生成元的计算是标准的半群导数。
\end{proof}

这一半群结构揭示了遗忘的\textbf{可加性}：两次连续的指数遗忘等价于一次具有
\textbf{乘积保留率}的遗忘。$n$次保留率为$\lambda$的遗忘等价于一次保留率为$\lambda^n$的遗忘——
遗忘在时间上的复合是指数加速的。这一观察为分层遗忘策略（多时间尺度的遗忘）提供了代数基础。

\subsection{遗忘算子视角下的相变定理}

遗忘算子理论为记忆-遗忘相变定理（定理~\ref{thm:phase_transition}）提供了
一个优雅的重新推导。在算子语言下：

\begin{enumerate}[label=(\arabic*)]
    \item 遗忘算子$\Fforget_\lambda$产生稳态记忆分布$\mu_\lambda$，
          其特征由\textbf{偏差泛函}$B(\lambda) = \norm{\E_{\mu_\lambda}[m] - \bar{\delta m}}_{\Mmem}$
          和\textbf{方差泛函}$V(\lambda) = \Tr(\Cov_{\mu_\lambda}(m))$刻画。
    \item 相变发生在$\frac{\dd}{\dd \lambda}(B(\lambda)^2 + V(\lambda))\big|_{\lambda=1}$改变符号时——
          即偏差随遗忘减少的边际收益超过方差随遗忘增加的边际代价时。
    \item 临界漂移$v_c$的特征方程：
          \begin{equation}
              \frac{\dd B(\lambda)^2}{\dd \lambda}\Big|_{\lambda=1}
              = -\frac{\dd V(\lambda)}{\dd \lambda}\Big|_{\lambda=1}。
          \end{equation}
          左侧是遗忘带来的偏差减少的边际收益，
          右侧是遗忘带来的方差增加的边际代价。
\end{enumerate}

这一``边际收益=边际代价''的结构在算子理论中是普遍的——
类似的方程出现在统计学习理论的偏差-方差权衡、
信息论中的率失真理论、以及热力学中的自由能最小化。

% ===========================================================================
% SECTION 7: DISCUSSION AND HONEST CRITIQUE
% ===========================================================================
\section{讨论与诚实暴击}

\subsection{主要贡献的回顾}

本文从SCX框架的Spring SE-1动力学出发，建立了记忆-遗忘的形式理论。
四个定理构成了一个逻辑递进的体系：

\begin{enumerate}[label=(\arabic*)]
    \item \textbf{遗忘的必要性定理}确立了遗忘的\textbf{逻辑地位}——不是工程优化，而是正确性条件。
    \item \textbf{记忆-遗忘相变定理}划分了\textbf{策略空间}——根据漂移速率，系统自动落入记忆最优或遗忘最优的区域。
    \item \textbf{最优遗忘界}刻画了\textbf{性能前沿}——遗忘与延迟的不可逃避的权衡，由信息论下界保证。
    \item \textbf{遗忘算子$\Fforget_t$}提供了\textbf{代数基础}——统一的算子语言，收敛性保证，以及与已知遗忘策略的精确对应。
\end{enumerate}

\subsection{诚实暴击：每个定理的局限与未竟问题}

以下是本文的\textbf{完整诚实暴击}——逐条列出每个结论的局限、简化假设和已知的反例或修正需求。

\subsubsection*{对遗忘必要性定理的暴击}

\begin{enumerate}[label=\textbf{暴击\,\arabic*:}, leftmargin=*]
    \item \textbf{渐近发散的有限时间无意义性。}
          定理证明$\lim_{t\to\infty} \E[\norm{\theta_t - \theta_t^*}^2] = \infty$，
          但任何实际系统运行在有限时间内。如果漂移$v_{\min}$足够小（如$10^{-6}$/步），
          发散在$10^9$步之后才变得显著——远超任何系统的运行寿命。此时``必然发散''
          在实际意义上等于``不发散''。需要有限时间的非渐近界来补充渐近结论。

    \item \textbf{漂移下界假设的可验证性问题。}
          $\liminf_{t \to \infty} \E[\norm{\theta_{t+1}^* - \theta_t^*}] \geq v_{\min} > 0$
          是定理的核心假设，但在实践中漂移本身是\textbf{不可直接观测}的——我们只能通过数据
          间接推断漂移的存在和大小。如果检测不到漂移（定理3的延迟问题），
          就无法判断定理1的前提是否成立。这是一个\textbf{认识论循环}。

    \item \textbf{均值回复漂移的例外。}
          如果漂移是均值回复的（如Ornstein-Uhlenbeck过程），历史平均$\bar{\theta}_t^*$
          可能是比当前$\theta_t^*$更好的估计（因为$\theta_t^*$受暂时性冲击的污染）。
          此时``遗忘''反而有害——我们需要更精细的\textbf{选择性遗忘}（仅遗忘永久性偏移，
          保留暂时性波动中的信息）。

    \item \textbf{单调记忆的夸大。}
          我们将``无遗忘''建模为等权平均$w(t, s) = 1/t$，但实际的Spring SE-1
          在递减步长$\alpha_t$下已经具有\textbf{隐式遗忘}——早期梯度因为$\alpha_s$大
          而产生大更新，但大$\alpha_s$时代已经过去（$\alpha_t \to 0$）。
          步长衰减本身构成了某种``遗忘''。因此，严格来说Spring SE-1并非完全不遗忘——
          它只是没有\textbf{显式}遗忘。
\end{enumerate}

\subsubsection*{对记忆-遗忘相变定理的暴击}

\begin{enumerate}[label=\textbf{暴击\,\arabic*:}, leftmargin=*]

    \item \textbf{PL条件的致命局限。}
          相变定理在PL条件全局成立的前提下推导了$v_c$的表达式。
          但PL条件在非凸深度学习中\textbf{几乎从不全局成立}——它只在最优点的某个
          吸引盆地内成立。当漂移$v$较大时，$\theta_t^*$可能在不同盆地向跳跃，
          PL条件在盆地边界处失效。此时``相变''可能被盆地跳跃的不连续性所掩盖。

    \item \textbf{固定步长的简化。}
          实际Spring SE-1使用递减步长$\alpha_t \propto 1/t$。
          在递减步长下，``有效遗忘率''是$\alpha_t$衰减率与显式遗忘率$\gamma$的
          组合效应。固定步长分析可能高估了遗忘的必要性（因为步长衰减本身提供了
          某种``渐进遗忘''）。

    \item \textbf{二次损失近似的精度。}
          偏差项推导$b(\lambda)^2 \approx v^2/(\alpha^2\mu^2) \cdot (\lambda/(1-\lambda))^2$
          使用了$g_t \approx H(\theta_t - \theta_t^*)$的线性化。在一般非凸损失下，
          梯度与参数误差的关系是非线性的，偏差的精确行为需要更复杂的分析。
          特别地，在损失景观具有鞍点和局部极小值的区域，线性近似可能完全失效。

    \item \textbf{常数$\lambda^*$的超参数敏感性。}
          最优保留率$\lambda^*$以复杂的方式依赖于问题参数$(\mu, L, \sigma^2, v)$——
          这些参数在实践中难以精确估计。$\lambda^*$对$\hat{v}$（漂移速率估计）的敏感性
          在临界区$(v \approx v_c)$可能极高——小的估计误差导致大的策略偏差。

\end{enumerate}

\subsubsection*{对最优遗忘界的暴击}

\begin{enumerate}[label=\textbf{暴击\,\arabic*:}, leftmargin=*]

    \item \textbf{高斯假设的普适性。}
          定理中使用的高斯观测模型是信息论界的``最有利情况''（least favorable case
          for detection within the class of distributions with bounded second moment
          is Gaussian）。在重尾分布或离散分布下，检测可能更难或更容易，
          下界常数的精确值会改变——但$\gamma \cdot \tau = \Omega(1/\delta^2)$的标度关系
          具有更广泛的普适性（由局部渐近正态性保证）。

    \item \textbf{偏移大小$\delta$的先验知识假设。}
          帕累托前沿的位置依赖于偏移大小$\delta$——但在实践中，$\delta$在偏移发生前
          \textbf{不可知}。这意味着不存在一个在所有可能的$\delta$下都帕累托最优的
          遗忘策略。设计者必须在\textbf{先验猜测}的$\delta$下优化，或使用极小极大准则。

    \item \textbf{假阳性与假阴性的不对称代价。}
          定理仅处理检测的统计显著性（$\alpha$水平），未纳入假阳性（在无偏移时误报偏移）
          和假阴性（偏移发生后未检测到）的不同决策代价。在许多应用中（如医疗诊断、
          金融风控），假阴性的代价远高于假阳性——这会使帕累托前沿向更低遗忘率方向倾斜。

\end{enumerate}

\subsubsection*{对遗忘算子$\Fforget_t$的暴击}

\begin{enumerate}[label=\textbf{暴击\,\arabic*:}, leftmargin=*]

    \item \textbf{Hilbert空间假设的范围。}
          记忆空间的Hilbert结构假设（完备内积空间）允许我们使用压缩映射原理和谱分析。
          但实际记忆系统（如Transformer的KV缓存、LSTM的隐藏状态和细胞状态）的结构
          可能不具有自然的Hilbert空间结构。在这些情况下，算子理论的结论需要被重新检验。

    \item \textbf{压缩性公理的过强性。}
          公理A1要求全局压缩性$\norm{\Fforget_t(m, \delta m) - \Fforget_t(m', \delta m)} \leq \rho_t \norm{m - m'}$。
          但实际遗忘机制可能仅在特定方向上压缩（如仅在``过时信息子空间''上压缩，
          在``稳定信息子空间''上保持）。方向性遗忘算子（如投影遗忘）不符合全局压缩公理，
          但可能在实际中更有用。

    \item \textbf{遗忘半群的离散性。}
          我们将遗忘建模为离散时间半群，但连续时间遗忘（遗忘微分方程
          $\frac{\dd m}{\dd t} = -\gamma(t) \cdot (m - \mathcal{T}(\delta m))$）
          可能提供更丰富的数学结构——特别是与随机微分方程和扩散过程的联系。

\end{enumerate}

\subsection{开放问题}

以下问题超出了本文的范围，但构成了Temporal SCX未来研究议程的核心：

\begin{enumerate}[label=\textbf{OP\,\arabic*}, leftmargin=*]

    \item \textbf{自适应遗忘的收敛理论。}
          当$\rho_t$依赖于$\norm{\delta m_t}$时（``惊讶驱动的遗忘''），
          遗忘算子序列$\{\Fforget_t\}$是非平稳和非线性的。
          证明其在一般条件下的收敛性（或发散条件）需要随机压缩系统理论。
          初步数值实验（未在本文中报告）显示：惊讶驱动的遗忘在阶跃偏移下优于固定速率遗忘，
          但在渐变偏移下可能出现震荡。

    \item \textbf{选择性遗忘的信息准则。}
          遗忘必要性定理说``必须遗忘''，但它没有指示\textbf{遗忘什么}。
          一种有前景的方向是使用条件互信息$I(\text{当前}; \text{历史} \mid \text{状态})$
          来决定保留/丢弃——当历史信息在当前状态下不提供关于目标的额外信息时，
          它可以被安全遗忘。这将是SCX的条件互信息准则$I(Y; P \mid S) > 0$
          在时间维度上的推广。

    \item \textbf{多时间尺度遗忘的分层结构。}
          遗忘半群的可加性（命题~\ref{prop:semigroup}）暗示了分层遗忘的可能性——
          不同时间尺度的遗忘算子可以嵌套。快速遗忘层处理瞬时噪声，
          慢速遗忘层保留长期趋势。这与神经科学中多时间尺度记忆系统
          （工作记忆/短期记忆/长期记忆）的层级结构形成有趣的对应。

    \item \textbf{遗忘与省定理的统一框架。}
          省定理和遗忘必要性定理分别处理了\textbf{无输入退化}和\textbf{输入漂移不适应}
          两种情况。一个统一的数学结构——``广义省定理''——将两者作为同一退化现象的
          两种极限情形（零输入速率 vs 非零输入漂移），仍有待建立。
          初步方向：将省定理中的遗忘率$\gamma$与遗忘必要性定理中的漂移$v$
          统一为一个``有效退化率''参数。

    \item \textbf{遗忘-记忆相变在神经网络中的实验验证。}
          定理2的相变预测可以在受控实验中检验：在具有已知漂移速率的合成数据流上
          训练线性网络或小型MLP，扫描遗忘率$\lambda$并测量稳态MSE，
          验证MSE曲线在$v < v_c$和$v > v_c$时的定性差异。初步理论预测在论文
          提交时尚未完成系统实验。

    \item \textbf{遗忘算子的量子推广。}
          遗忘算子的压缩性（公理A1）与量子信道（完全正压缩映射）的数学形式高度相似。
          记忆状态的量子表示可能产生信息论上更优的遗忘策略——例如利用量子叠加
          同时保留``遗忘''和``不遗忘''的幅度，直到分布偏移被确认后再进行``测量''。

\end{enumerate}

\subsection{对SCX理论体系的影响}

Temporal SCX的四个定理对SCX框架的已有定理体系产生了以下深远影响：

\begin{enumerate}[label=(\arabic*)]
    \item \textbf{省定理的补全。}省定理处理了``停止一切更新''的退化场景。
          遗忘必要性定理处理了``持续更新但目标漂移''的退化场景。
          两者合在一起，覆盖了SCX系统可能面临的全部退化模式：系统内部的衰退（省定理）
          和系统外部的变化（遗忘必要性定理）。
    \item \textbf{Spring SE-1的推广。}将Spring SE-1从\textbf{收敛到固定目标}推广为
          \textbf{跟踪时变目标}，大大扩展了Spring动力学在真实世界场景中的适用范围。
    \item \textbf{SCX记忆库设计的理论基础。}Spring记忆库$\mathcal{M}_t$的设计
          现在有了一个形式化的指导原则：遗忘算子的选择（$\lambda$或$W$）应由
          环境的漂移速率$v$决定——在低漂移环境中扩展记忆，在高漂移环境中收缩记忆。
    \item \textbf{质变点定理的补充。}质变点定理描述了\textbf{学习过程的相变}
          （探索→收敛），记忆-遗忘相变定理描述了\textbf{记忆策略的相变}
          （记忆→遗忘）。两个相变共同描绘了SCX系统的完整动力学景观。
\end{enumerate}

\subsection{工程建议（带诚实标注）}

基于本文的定理，我们提出以下SCX系统设计建议——每个建议均标注了
从数学定理到工程实践的映射不确定性。

\begin{enumerate}[label=\textbf{建议\,\arabic*:}, leftmargin=*]

    \item \textbf{漂移速率估计作为一等公民。}
          \heuristic{} 定理1--3均表明：遗忘策略的选择应取决于漂移速率$v$。
          因此，SCX系统应内建漂移速率估计模块（如基于滑动窗口的$\theta_t^*$变化检测），
          并将估计值$\hat{v}$作为遗忘超参数的自动调节信号。
          \textbf{不确定性：}漂移速率估计本身需要遗忘（旧数据中的$\theta_s^*$已经过时）——
          这构成了另一个层的``元遗忘''问题。

    \item \textbf{分层遗忘架构。}
          \heuristic{} 在半群可加性的指导下，不同组件可以使用不同的保留率：
          Spring参数更新使用较慢的遗忘（$\lambda \approx 0.99$），
          而Yajie多专家审计使用较快的遗忘（$\lambda \approx 0.9$），
          以快速检测数据质量的变化。
          \textbf{不确定性：}组件之间的遗忘率差异可能导致不一致的内部状态表示。

    \item \textbf{遗忘-延迟预算。}
          \heuristic{} 定理3的帕累托前沿将$\gamma \cdot \tau$视为不可突破的预算约束。
          设计者应根据应用的延迟容忍度（最大可接受$\tau$）来确定最大可接受$\gamma$，
          然后在$\gamma \leq \gamma_{\max}$的约束下优化其他指标。
          \textbf{不确定性：}帕累托界是渐近的——在有限样本下，存在以更高方差为代价的
          ``超前沿''策略（如序贯概率比检验SPRT），其有限样本性能可能突破渐近界。

\end{enumerate}

% ===========================================================================
% SECTION 8: CONCLUSION
% ===========================================================================
\section{结论：记忆的数学本质}

本文的核心论点是：\textbf{记忆与遗忘不是认知心理学的研究对象——它们在数学上具有
严格的、可从SCX框架的动力学中推导出来的形式结构。}
这一形式结构揭示了关于记忆的三个本质事实：

\begin{enumerate}[label=(\arabic*)]
    \item \textbf{遗忘是逻辑必然（定理1）。}在非稳态数据流中，无遗忘系统必然发散——
          这不是工程权衡，而是数学定理。遗忘不是记忆系统的缺陷，
          而是维持有限跟踪误差的\textbf{必要条件}。
    \item \textbf{记忆-遗忘存在相变（定理2）。}环境漂移速率穿越临界值时，
          最优策略在``记忆一切''和``遗忘大部分''之间发生定性跃迁。
          这一相变在结构上平行于学习过程中的质变点，暗示两者可能有共同的深层数学根源。
    \item \textbf{遗忘的速度存在不可突破的下界（定理3）。}遗忘与检测延迟的乘积
          受信息论下界约束——不存在同时快速遗忘和快速检测的方案。
          这是\textbf{时间版本的测不准原理}：遗忘速度和检测确定性的乘积存在不可突破的下界。
\end{enumerate}

\textbf{终极隐喻：}记忆是一条河流——如果水不流动（无遗忘），它会变成死水（发散）；
如果水流太快（过度遗忘），它无法反映两岸的风景（高方差）。
最优的记忆系统像一条稳态的河流：水流的速度（遗忘率）恰好匹配
河床的变化速率（漂移速率），在流动中保持形态的恒定性。
这一隐喻并非文学修辞——它是定理2的平衡条件$\frac{\dd\MSE}{\dd\lambda} = 0$
在物理直觉中的投射。

\metaphor{} \textbf{类比标注：}终极隐喻是启发式的物理类比，不具有数学效力。
河流-记忆的类比在以下意义上成立：两者均涉及``流入-流出平衡下的稳态形态''，
且稳态由系统参数（水流速度/遗忘率，河床变化/漂移速率）唯一决定。
然而，河流是连续介质力学的对象，记忆是信息论的对象——
两者在数学结构上的精确对应（如果存在）仍有待建立。

\vspace{12pt}
\begin{center}
\rule{0.5\textwidth}{0.4pt}
\end{center}
\vspace{12pt}

\textbf{致谢：}感谢SCX理论体系（2026）提供了Spring SE-1、省定理、质变点定理的数学基础，
本文的所有推广均建立在这些已有成果之上。所有错误和过度简化由本文作者负责。

% ===========================================================================
% APPENDIX
% ===========================================================================
\appendix

\section{有限窗口遗忘的相变分析}

本附录补充定理~\ref{thm:phase_transition}在有限窗口遗忘策略下的对应分析。

\begin{proposition}[有限窗口遗忘的临界窗口]\label{prop:fw_critical}
在有限窗口遗忘策略$w(t, s) = 1/W \cdot \mathbf{1}[t-s \leq W]$下，
渐近MSE为：
\begin{equation}
    \MSE(W) = \frac{v^2 W^2}{4\alpha^2\mu^2} + \frac{\alpha\sigma^2}{\mu W}。
\end{equation}

最优窗口大小$W^*$由一阶条件给出：
\begin{equation}
    W^* = \left(\frac{2\alpha^3\sigma^2\mu}{v^2}\right)^{1/3}。
\end{equation}

当$v \to 0$时，$W^* \to \infty$（无限窗口，即无遗忘）；
当$v \to \infty$时，$W^* \to 0$（即时遗忘）。
临界漂移（定义为$W^* = 1$时的$v$值）为$v_c^{\text{FW}} = \sqrt{2\alpha^3\sigma^2\mu}$。
\end{proposition}

\begin{proof}
与定理~\ref{thm:phase_transition}的证明平行。在窗口大小$W$下：
\begin{itemize}
    \item 有效偏差：历史平均滞后$\frac{1}{W}\sum_{k=0}^{W-1}(\theta_{t-k}^* - \theta_t^*)$。
          在线性漂移下，$\norm{\frac{1}{W}\sum_{k=0}^{W-1}(\theta_{t-k}^* - \theta_t^*)}^2 = \frac{v^2 W^2}{4}$。
    \item 有效方差：$W$个独立样本的平均的方差为$\alpha^2\sigma^2/W$。
\end{itemize}
求解$\frac{\dd \MSE(W)}{\dd W} = 0$得到$W^*$。
\end{proof}

指数遗忘和有限窗口遗忘在以下极限下等价：
当$\lambda = 1 - 1/W$且$W$大时，$N_{\text{eff}}^{\text{EF}} \approx 2W$，
而$N_{\text{eff}}^{\text{FW}} = W$。因此，在有效样本量的意义上，
指数遗忘的``有效窗口''是有限窗口的两倍（因为它对近期历史给予更高权重）。

\section{遗忘算子的谱分析}

\begin{proposition}[遗忘算子的谱]\label{prop:spectrum}
在有限维记忆空间$\Mmem \cong \R^d$中，标准遗忘算子
$\Fforget_\lambda(m) = \lambda m + (1-\lambda)\mathcal{T}(\delta m)$
的谱满足：
\begin{equation}
    \sigma(\Fforget_\lambda) = \{\lambda\} \cup \{\lambda \cdot \sigma_i : \sigma_i \in \sigma(P_{\Mmem}\mathcal{T})\}，
\end{equation}
其中$P_{\Mmem}$是到$\Mmem$的正交投影。特别地，遗忘算子的谱半径
$\rho(\Fforget_\lambda) = \lambda$（由公理A1的压缩性保证）。
\end{proposition}

\begin{proof}
在$\R^d$中，$\Fforget_\lambda$可表示为矩阵：
\begin{equation}
    F_\lambda = \lambda I_d + (1-\lambda)T，
\end{equation}
其中$T$是$\mathcal{T}$在$\Mmem$的某组基下的矩阵表示。$F_\lambda$的特征值为
$\lambda + (1-\lambda)\tau_i$，其中$\tau_i$是$T$的特征值。
由于$\norm{T} \leq 1$（非膨胀性），$\abs{\tau_i} \leq 1$，因此$\abs{\lambda + (1-\lambda)\tau_i} \leq 1$。
谱半径$\max_i \abs{\lambda + (1-\lambda)\tau_i} = \lambda + (1-\lambda) = 1$当$\tau_i = 1$时达到，
或$\lambda$当所有$\abs{\tau_i} < 1$时。
\end{proof}

谱分析揭示了遗忘算子的\textbf{混合性质}：在$\tau_i = 1$的方向上（对应$\mathcal{T}$的
不动点子空间），$\Fforget_\lambda$的特征值为1（该方向无遗忘——信息被完美保留）；
在$\tau_i = 0$的方向上，特征值为$\lambda$（该方向以速率$\lambda$被遗忘）。
这一方向性遗忘结构超出了公理A1的简单压缩模型——
它暗示了\textbf{结构化遗忘}的可能性：在不同的信息子空间上以不同的速率遗忘。

\openproblem{} \textbf{开放问题：}结构化遗忘（方向依赖性遗忘率的遗忘算子）
的收敛性分析及其在深度学习中的应用。初步猜想：Transformer注意力机制中的
``注意力衰减''（attention dropout / stochastic depth）可被理解为
在高维记忆空间的特定子空间上的结构化遗忘。

% ===========================================================================
% BIBLIOGRAPHY
% ===========================================================================
\begin{thebibliography}{99}

\bibitem{scx2026theorems}
SCX Theory Group.
\emph{SCX Core Theorems: Economy Theorem, Qualitative Change Point Theorem, and Multi-Expert Consistency}.
SCX Technical Report, 2026.

\bibitem{scx2026spring_convergence}
SCX Theory Group.
\emph{Spring SE-1 Convergence Analysis: Two-Phase Dynamics of Robbins-Monro SGD on State Space}.
SCX Technical Report, 2026.

\bibitem{scx2026cc_audit}
SCX Theory Group.
\emph{SCX Framework for Label Noise Detection: The Five-Layer Architecture}.
SCX Technical Report, 2026.

\bibitem{robbins1951}
H.~Robbins and S.~Monro.
\emph{A Stochastic Approximation Method}.
Annals of Mathematical Statistics, 22(3):400--407, 1951.

\bibitem{polyak1963}
B.~T.~Polyak.
\emph{Gradient methods for the minimisation of functionals}.
USSR Computational Mathematics and Mathematical Physics, 3(4):864--878, 1963.

\bibitem{lojasiewicz1963}
S.~Łojasiewicz.
\emph{A topological property of real analytic subsets}.
Coll. du CNRS, Les équations aux dérivées partielles, 117:87--89, 1963.

\bibitem{karimi2016}
H.~Karimi, J.~Nutini, and M.~Schmidt.
\emph{Linear convergence of gradient and proximal-gradient methods under the Polyak-Łojasiewicz condition}.
ECML-PKDD, 2016.

\bibitem{chernoff1952}
H.~Chernoff.
\emph{A measure of asymptotic efficiency for tests of a hypothesis based on the sum of observations}.
Annals of Mathematical Statistics, 23(4):493--507, 1952.

\bibitem{stein1952}
C.~Stein.
\emph{On the theory of testing composite hypotheses}.
Annals of Mathematical Statistics, 23:148--149, 1952.

\bibitem{lecun1986}
Y.~Le~Cam.
\emph{Asymptotic Methods in Statistical Decision Theory}.
Springer, 1986.

\bibitem{hoeffding1963}
W.~Hoeffding.
\emph{Probability inequalities for sums of bounded random variables}.
JASA, 58(301):13--30, 1963.

\bibitem{ebbinghaus1885}
H.~Ebbinghaus.
\emph{Über das Gedächtnis: Untersuchungen zur experimentellen Psychologie}.
Duncker \& Humblot, 1885.

\bibitem{wixted2007}
J.~T.~Wixted and S.~K.~Carpenter.
\emph{The Wickelgren Power Law and the Ebbinghaus forgetting curve}.
In H.~L.~Roediger, III (Ed.), Cognitive psychology of memory, 2007.

\bibitem{gama2014}
J.~Gama, I.~Žliobaitė, A.~Bifet, M.~Pechenizkiy, and A.~Bouchachia.
\emph{A survey on concept drift adaptation}.
ACM Computing Surveys, 46(4):1--37, 2014.

\bibitem{zinkevich2003}
M.~Zinkevich.
\emph{Online convex programming and generalized infinitesimal gradient ascent}.
ICML, 2003.

\bibitem{hazan2016}
E.~Hazan.
\emph{Introduction to Online Convex Optimization}.
Foundations and Trends in Optimization, 2016.

\bibitem{nagel2022}
R.~Nagel (Ed.).
\emph{One-Parameter Semigroups of Positive Operators}.
Springer Lecture Notes in Mathematics, 1986/2022.

\end{thebibliography}

\end{document}
